\section{秋季星空與天體}
\begin{itemize}
    \item 秋季四邊形：飛馬座\(\alpha\)（室宿一）、飛馬座\(\beta\)（室宿二）、飛馬座\(\gamma\)（壁宿一）、仙女座\(\alpha\)（壁宿二）
    \item 秋季南三角：鯨魚座的土司空、南魚座的北落師門、鳳凰座的火鳥六
    \item 秋季認星歌訣：\\秋夜北斗靠地平；仙后五星空中昇；仙女一字指東北；飛馬凌空四邊形\\英仙星座照夜空；大陵五星光會變；南天寂靜亮星少；北落師門賽明燈
\end{itemize}

\subsection{飛馬座 Pegasus}
\begin{enumerate}
    \item 秋季四邊形：飛馬座\(\alpha\)（室宿一）、飛馬座\(\beta\)（室宿二）、飛馬座\(\gamma\)（壁宿一）
    \item 面積第七
    \item 危宿三（飛馬座\(\epsilon\)）:遠距雙星，2.4 等黃色超巨星及一顆 8 等星。代表飛馬座的鼻子。
    \item 室宿二（飛馬座\(\beta\)）：紅巨星\&不規則變星，星等變化在2.2到2.75間
\end{enumerate}
盾牌座中的深空天體：
\newline

\begin{center}
\begin{longtable}{c | c | c | c | c | c}
    % \hline
    編號 & 種類 & 別名 & 視尺寸 & RA & Dec \\
    \hline
    \endfirsthead

    % \hline
    % 編號 & 種類 & 別名 & 視尺寸 & 拍攝技巧 & 曝光時間 \\
    % \hline
    % \endhead

    % \hline
    % \multicolumn{6}{r}{接下頁} \\
    \endfoot

    % \hline
    \endlastfoot

    M15/NGC 7078 & 球狀星團 & - & \(\ang{;18;}\) & \(21^\text{h }29^\text{m }57^\text{s}\) & \(+\ang{12;10;0}\) \\ %RGB & - \\
    NGC 7331 & 無棒螺旋星系 & - & \(\ang{;10;30}\times\ang{;7;42}\) & \(22^\text{h }37^\text{m }041^\text{s}\) & \(+\ang{34;24;55}\) \\ %RGB & - \\
    \makecell{NGC 7317\\NGC 7318a\\NGC 7318b\\NGC 7319\\NGC 7320c} & - & 史蒂芬五重奏星系 & \(\ang{;6;}\) & \(22^\text{h }35^\text{m }58^\text{s}\) & \(+\ang{33;57;36}\) \\ %RGB/LRGBH\(\alpha\) & - \\
    NGC 7479 & 棒旋星系 & - & \(\ang{;4;6}\times\ang{;3;6}\) & \(23^\text{h }04^\text{m }56^\text{s}\) & \(+\ang{12;19;22}\) \\ %RGB & - \\
    NGC 7814 & 螺旋星系 & - & \(\ang{;5;30}\times\ang{;2;18}\) & \(0^\text{h }03^\text{m }15^\text{s}\) & \(+\ang{16;8;43}\) \\ %RGB & - \\
\end{longtable}
\end{center}

\subsection{仙女座 Andromeda}
\begin{enumerate}
    \item 最亮星：壁宿二（仙女座\(\alpha\)）視星等2.1，曾為飛馬座的一部分。是一對聯星。象徵仙女安朵美達的頭，但阿拉伯傳統星名意為「駿馬的肚臍」
    \item 奎宿九（仙女座\(\beta\)）：是紅巨星，視星等為2.06
    \item 天大將軍一（仙女座\(\gamma\)）：為雙星系統，綜合視星等為2.14\\主星（\(\gamma^1\)）：金黃色，視星等2.3\\伴星（\(\gamma^2\)）：靛藍色，視星等5.0
    \item 室宿一和對角的壁宿二連線延伸\(\rightarrow\)紅色亮星（奎宿九）\(\rightarrow\)橘色亮星（天大將軍一）\\\(\Rightarrow\)找到仙女座
\end{enumerate}
仙女座中的深空天體：
\newline

\begin{center}
\begin{longtable}{c | c | c | c | c | c}
    % \hline
    編號 & 種類 & 別名 & 視尺寸 & RA & Dec \\
    \hline
    \endfirsthead

    % \hline
    % 編號 & 種類 & 別名 & 視尺寸 & 拍攝技巧 & 曝光時間 \\
    % \hline
    % \endhead

    % \hline
    % \multicolumn{6}{r}{接下頁} \\
    \endfoot

    % \hline
    \endlastfoot

    M31/NGC 224 & 棒旋星系 & 仙女座大星系 & \(\ang{3;10;}\times\ang{1;;}\) & \(0^\text{h }42^\text{m }43^\text{s}\) & \(+\ang{41;16;18}\) \\ %RGB/LRGBH\(\alpha\) & > 2 hr \\
    M32/NGC 221 & 橢圓星系 & （M31伴星系） & \(\ang{;7;42}\) & \(0^\text{h }42^\text{m }42^\text{s}\) & \(+\ang{40;51;55}\) \\ %RGB & - \\
    M110/NGC 205 & 橢圓星系 & （M31伴星系） & \(\ang{;21;54}\times\ang{;11;}\) & \(0^\text{h }40^\text{m }22^\text{s}\) & \(+\ang{41;41;07}\) \\ %RGB & - \\
    C22/NGC 7662 & 行星狀星雲 & 藍雪球星雲 & \(\ang{;;32}\times\ang{;;28}\) & \(23^\text{h }25^\text{m }24^\text{s}\) & \(+\ang{42;32;05}\) \\ %RGB/LRGBH\(\alpha\)/HOO & - \\
    C28/NGC 752 & 疏散星團 & - & \(\ang{;50;}\) & \(01^\text{h }57^\text{m }34^\text{s}\) & \(+\ang{37;50;00}\) \\ %RGB & - \\
    C27/NGC 891 & 棒旋星系 & - & \(\ang{;7;24}\) & \(02^\text{h }22^\text{m }32^\text{s}\) & \(+\ang{42;20;56}\) \\ %RGB & - \\
\end{longtable}
\end{center}

\subsection{仙王座 Cepheus}
\begin{enumerate}
    \item 形狀成一鉛筆頭，因為位於天球北極，因此一年四季皆看得到
    \item 最亮星：仙王座\(\alpha\)（天鉤五）：視星等2.45
    \item 仙王座\(\delta\)（造父一）：雙星及變星，亮度變化範圍從 3.5 等到 4.4 等（較亮的是顆黃巨星），週期為\(5\frac{3}{8}\)天
    \item 仙王座\(\beta\)：雙星及變星，3.2等的藍巨星與7.9 等的伴星
    \item 仙王座\(\mu\) （柘榴星）：變星，亮度變化範圍從3.4等到5.1等，週期為兩年。紅巨星，顏色深紅而得柘榴星之名
\end{enumerate}
仙王座中的深空天體：
\newline

\begin{center}
\begin{longtable}{c | c | c | c | c | c}
    % \hline
    編號 & 種類 & 別名 & 視尺寸 & RA & Dec \\
    \hline
    \endfirsthead

    % \hline
    % 編號 & 種類 & 別名 & 視尺寸 & 拍攝技巧 & 曝光時間 \\
    % \hline
    % \endhead

    % \hline
    % \multicolumn{6}{r}{接下頁} \\
    \endfoot

    % \hline
    \endlastfoot

    C4/NGC 7023 & 反射/暗星雲 & 鳶尾花星雲 & 暗星雲 > \(\ang{;18;}\) & \(21^\text{h }01^\text{m }36^\text{s}\) & \(+\ang{68;10;10}\) \\ %RGB & > 10 hr \\
    C9/Sh2-155 & 發射/反射/暗星雲 & 洞穴星雲 & \(\ang{;50;}\times\ang{;30;}\) & \(22^\text{h }57^\text{m }17^\text{s}\) & \(+\ang{62;28;33}\) \\ %皆可 & > 6 hr \\
    Sh2-136 & 反射星雲/包克雲球 & Ghost Nebula & \(\ang{;30;}\) & \(21^\text{h }16^\text{m }26^\text{s}\) & \(+\ang{68;15;36}\) \\ %RGB & > 10 hr \\
    NGC 6946 & 中間螺旋星系 & 焰火星系 & \(\ang{;16;}\times\ang{;11;12}\) & \(20^\text{h }34^\text{m }51^\text{s}\) & \(+\ang{60;9;14}\) \\ %RGB/LRGBH\(\alpha\) & - \\
    IC 1396 & 發射星雲 & \makecell{象鼻星雲\\（IC 1396A）} & \(\ang{2;50;}\) & \(21^\text{h }38^\text{m }57^\text{s}\) & \(+\ang{57;29;20}\) \\ %RGB/SHO & > 4 hr \\
    LDN 1235 & 反射/暗星雲 & 鯊魚星雲 & \(\ang{2;;}\times\ang{1;30;}\) & \(22^\text{h }14^\text{m }51^\text{s}\) & \(\ang{73;32;48}\) \\ %RGB & > 10 hr \\
    NGC 7380 & 發射星雲/疏散星團 & 巫師星雲 & \(\ang{;25;}\) & \(22^\text{h }47^\text{m }21^\text{s}\) & \(+\ang{58;07;56}\) \\ %RGB/SHO & > 6 hr \\
    Barnard 150 & 暗星雲 & 海馬星雲 & \(\ang{1;;}\) & \(20^\text{h }50^\text{m }40^\text{s}\) & \(+\ang{60;18;0}\) \\
    Barnard 175 & 包克雲球 & - & \(\ang{1;;}\) & \(22^\text{h }13^\text{m }33^\text{s}\) & \(+\ang{70;15;25}\) \\
    Sh2-132 & 發射星雲 & - & \(\ang{1;30}\times\ang{1;30;}\) & \(22^\text{h }19^\text{m }09^\text{s}\) & \(+\ang{56;04;45}\)
\end{longtable}
\end{center}

\subsection{仙后座 Cassiopeia}
\begin{enumerate}
    \item 五顆亮星呈W形可用於辨認北極星。（W的交點延長五倍可找到北極星）
    \item 位於天球北極，因此一年四季皆看得到
\end{enumerate}
仙后座中的深空天體：
\newline

\begin{center}
\begin{longtable}{c | c | c | c | c | c}
    % \hline
    編號 & 種類 & 別名 & 視尺寸 & RA & Dec \\
    \hline
    \endfirsthead

    % \hline
    % 編號 & 種類 & 別名 & 視尺寸 & 拍攝技巧 & 曝光時間 \\
    % \hline
    % \endhead

    % \hline
    % \multicolumn{6}{r}{接下頁} \\
    \endfoot

    % \hline
    \endlastfoot

    M52/NGC 7654 & 疏散星團 & - & \(\ang{;13;}\) & \(23^\text{h }24^\text{m }47^\text{s}\) & \(+\ang{61;35;35}\) \\ %RGB & - \\
    M103/NGC 581 & 疏散星團 & - & \(\ang{;6;}\) & \(01^\text{h }33^\text{m }21^\text{s}\) & \(+\ang{60;39;28}\) \\ %RGB & - \\
    NGC 457 & 疏散星團 & 貓頭鷹星團 & \(\ang{;20;}\) & \(01^\text{h }19^\text{m }32^\text{s}\) & \(+\ang{58;17;26}\) \\ %RGB & - \\
    NGC 663 & 疏散星團 & - & \(\ang{;16;}\) & \(01^\text{h }46^\text{m }15^\text{s}\) & \(+\ang{61;13;05}\) \\ %RGB & - \\
    NGC 281/IC 11 & 發射星雲 & 小精靈星雲 & \(\ang{;35;}\) & \(0^\text{h }52^\text{m }59^\text{s}\) & \(+\ang{56;37;18}\) \\ %RGB/SHO & > 4 hr \\
    C11/NGC 7635 & 發射星雲 & 氣泡星雲 & \(\ang{;15;}\times\ang{;8;}\) & \(23^\text{h }20^\text{m }45^\text{s}\) & \(+\ang{61;12;44}\) \\ %RGB/SHO/HOO & > 8 hr \\
    IC 1805/Sh2-190 & 發射星雲 & 心臟星雲 & \(\ang{2;30;}\) & \(2^\text{h }32^\text{m }41^\text{s}\) & \(+\ang{61;27;24}\) \\ %RGB/SHO & > 4 hr \\
    IC 1795/NGC 896 & 發射星雲 & \makecell{魚頭星雲\\（屬於心臟）} & \(\ang{;30;}\) & \(2^\text{h }25^\text{m }27^\text{s}\) & \(+\ang{62;1;9}\) \\ %RGB/SHO & - \\
    W5/Sh2-199/IC 1848 & 發射星雲 & 靈魂星雲 & \(\ang{2;30;}\times\ang{1;15;}\) & \(02^\text{h }51^\text{m }11^\text{s}\) & \(+\ang{60;24;8}\) \\ %RGB/SHO & > 4 hr \\
    Sh2-185 & 發射/反射星雲 & \makecell{Glost of\\Cassiopeia} & \(\ang{2;;}\) & \(01^\text{h }00^\text{m }00^\text{s}\) & \(+\ang{60;59;00}\)
\end{longtable}
\end{center}

\subsection{英仙座 Perseus}
\begin{enumerate}
    \item 最亮星：天船三，一顆黃白超巨星，和周圍的恆星組成了有名的英仙座\(\alpha\)星團
    \item 英仙座\(\beta\)星（大陵五）又稱梅杜莎之眼
    \item 最早被人類所認知的食變星
    \item 光變周期2.867日 星等在+2.1到+3.3間變化
    \item 形象：為一手持寶劍，一手抓著蛇髮女妖梅杜莎的首級
    \item 英仙座流星雨：英仙座流星雨是北半球三大流星雨之一，也是著名的七大流星雨之一，英仙座流星雨的數量多，他被稱為是最適合非專業流星觀測者的流星雨
\end{enumerate}
英仙座中的深空天體：
\newline

\begin{center}
\begin{longtable}{c | c | c | c | c | c}
    % \hline
    編號 & 種類 & 別名 & 視尺寸 & RA & Dec \\
    \hline
    \endfirsthead

    % \hline
    % 編號 & 種類 & 別名 & 視尺寸 & 拍攝技巧 & 曝光時間 \\
    % \hline
    % \endhead

    % \hline
    % \multicolumn{6}{r}{接下頁} \\
    \endfoot

    % \hline
    \endlastfoot

    M34/NGC 1039 & 疏散星團 & - & \(\ang{;35;}\) & \(2^\text{h }42^\text{m }06^\text{s}\) & \(+\ang{42;44;46}\) \\ %RGB & - \\
    NGC 869 + NGC 884 & \makecell{雙星團\\（疏散星團）} & - & \(\ang{1;;}\times\ang{;30;}\) & \(2^\text{h }18^\text{m }59^\text{s}\) & \(+\ang{57;07;02}\) \\ %RGB & - \\
    M76 & 行星狀星雲 & 小啞鈴星雲 & \(\ang{;2;42}\times\ang{;1;48}\) & \(01^\text{h }42^\text{m }20^\text{s}\) & \(+\ang{51;34;31}\) \\ %RGB/HOO & - \\
    NGC 1499 & 發射星雲 & 加州星雲 & \(\ang{2;30}\times\ang{1;;}\) & \(4^\text{h }03^\text{m }14^\text{s}\) & \(+\ang{36;22;2}\) \\ %RGB/HOO/SHO & > 4 hr \\
    NGC 1333 & 廣域暗星雲 & 英仙座分子雲 & - & \(03^\text{h }28^\text{m }54^\text{s}\) & \(+\ang{31;22;12}\) \\ %
\end{longtable}
\end{center}

\subsection{蝎虎座 Lacerta}
\begin{enumerate}
    \item 最亮星：蝎虎座\(\alpha\)星（螣蛇一），3.8 等
    \item 位於銀河邊緣
    \item 拉丁文意思是蜥蜴
    \item 和仙后座一樣呈W形
    \item 曾與仙女座內的幾顆恆星歸為腓特烈榮譽座，以紀念普魯士腓特烈大帝，但不久後便不再通用
\end{enumerate}

\subsection{小馬座 Equuleus}
\begin{enumerate}
    \item 最亮星：小馬座\(\alpha\)星（虛宿）黃色巨星，視星等3.9等
    \item 全天第二小星座
\end{enumerate}

\subsection{鯨魚座 Cetus}
\begin{enumerate}
    \item 最亮星：鯨魚座\(\beta\)土司空，視星等2.04
    \item 形象：希臘人認為牠是海生怪物，牠擁有一個血盆大口，前足如陸上生物，有一條蛇狀的尾
\end{enumerate}

\subsection{白羊座 Aries}
\begin{enumerate}
    \item 僅四顆高亮恆星：婁宿三、婁宿一、婁宿二和胃宿三
    \item 最亮星：婁宿三，視星等2.01
    \item 白羊座\(\gamma\)星（婁宿二）：雙星，兩顆4.6等藍白色恆星
    \item 白晝白羊座流星雨：從每年5月22日起持續到7月2日，是日間比較強烈的流星雨
\end{enumerate}
白羊座中的深空天體：

\begin{center}
\begin{longtable}{c | c | c | c | c | c}
    % \hline
    編號 & 種類 & 別名 & 視尺寸 & RA & Dec \\
    \hline
    \endfirsthead

    % \hline
    % 編號 & 種類 & 別名 & 視尺寸 & 拍攝技巧 & 曝光時間 \\
    % \hline
    % \endhead

    % \hline
    % \multicolumn{6}{r}{接下頁} \\
    \endfoot

    % \hline
    \endlastfoot

    LBN 762 + LBN 753 & 暗/反射星雲 & - & \(\ang{2;;}\text{ to }\ang{6;;}\) & \(3^\text{h }00^\text{m }01^\text{s}\) & \(+\ang{17;12;2}\) \\ % RGB & > 10 hr \\
\end{longtable}
\end{center}

\subsection{三角座 Triangulum}
\begin{enumerate}
    \item 最亮星：三角座\(\beta\)星（天大將軍九），白巨星，視星等3.00
    \item 婁宿增六：黃白次巨星，視星等3.41
    \item 三角座六：雙星，5 等及 7 等
    \item 因三顆較為明亮的恆星組成狹長三角形得名
\end{enumerate}
三角座中的深空天體：
\newline

\begin{center}
\begin{longtable}{c | c | c | c | c | c}
    % \hline
    編號 & 種類 & 別名 & 視尺寸 & RA & Dec \\
    \hline
    \endfirsthead

    % \hline
    % 編號 & 種類 & 別名 & 視尺寸 & 拍攝技巧 & 曝光時間 \\
    % \hline
    % \endhead

    % \hline
    % \multicolumn{6}{r}{接下頁} \\
    \endfoot

    % \hline
    \endlastfoot

    M33/NGC 598 & 螺旋星系 & 三角座星系 & \(\ang{1;10;48}\times\ang{;41;42}\) & \(01^\text{h }33^\text{m }51^\text{s}\) & \(+\ang{30;39;36}\) \\ %RGB/LRGBH\(\alpha\) & > 2 hr \\
\end{longtable}
\end{center}

\subsection{寶瓶座 Aquarius}
\begin{enumerate}
    \item 面積大但黯淡，面積排名第十
    \item 最亮星：虛宿一（寶瓶座\(\beta\)），視星等為2.90
    \item 寶瓶座\(\eta\)流星雨、寶瓶座\(\delta\)流星雨和寶瓶座\(\iota\)流星雨
\end{enumerate}
寶瓶座中的深空天體：
\newline

\begin{center}
\begin{longtable}{c | c | c | c | c | c}
    % \hline
    編號 & 種類 & 別名 & 視尺寸 & RA & Dec \\
    \hline
    \endfirsthead

    % \hline
    % 編號 & 種類 & 別名 & 視尺寸 & 拍攝技巧 & 曝光時間 \\
    % \hline
    % \endhead

    % \hline
    % \multicolumn{6}{r}{接下頁} \\
    \endfoot

    % \hline
    \endlastfoot

    M2/NGC7089 & 球狀星團 & - & \(\ang{;16;}\) & \(21^\text{h }33^\text{m }27^\text{s}\) & \(-\ang{0;49;27}\) \\ %RGB & - \\
    C55/NGC 7009 & 行星狀星雲 & 土星星雲 & \(\ang{;41;}\times\ang{;35;}\) & \(21^\text{h }04^\text{m }10^\text{s}\) & \(-\ang{11;21;47}\) \\ %RGB/HOO & - \\
    C63/NGC 7293 & 行星狀星雲 & 螺旋星雲/蝸牛星雲 & \(\ang{;25;}\) & \(22^\text{h }29^\text{m }39^\text{s}\) & \(-\ang{20;50;13}\) \\ %皆可 & > 6 hr \\
    M73/NGC 6994 & 星群 & - & \(\ang{;2;48}\) & \(20^\text{h }58^\text{m }55^\text{s}\) & \(-\ang{12;38;08}\) \\ %RGB & - \\
    M72/NGC 6981 & 球狀星團 & - & \(\ang{;6;36}\) & \(20^\text{h }53^\text{m }27^\text{s}\) & \(-\ang{12;32;12}\) \\ %RGB & - \\
\end{longtable}
\end{center}

\subsection{南魚座 Piscis Austrinus}
\begin{enumerate}
    \item 順著水瓶座倒的酒往下看即可找到
    \item 唯一的亮星：南魚座\(\alpha\)星（北落師門）是顆藍白色的巨星，全天第17亮星（不包括太陽），視星等為1.16
    \item 南魚座流星雨
\end{enumerate}

\subsection{雙魚座 Pisces}
\begin{enumerate}
    \item 雖然是較大的星座，但組成的恆星都很暗，容易辨別的特徵為兩個雙魚座小環
    \item 目前為春分點的星座
    \item 南雙魚座流星雨、北雙魚座流星雨
\end{enumerate}
雙魚座中的深空天體：
\newline

\begin{center}
\begin{longtable}{c | c | c | c | c | c}
    % \hline
    編號 & 種類 & 別名 & 視尺寸 & RA & Dec \\
    \hline
    \endfirsthead

    % \hline
    % 編號 & 種類 & 別名 & 視尺寸 & 拍攝技巧 & 曝光時間 \\
    % \hline
    % \endhead

    % \hline
    % \multicolumn{6}{r}{接下頁} \\
    \endfoot

    % \hline
    \endlastfoot

    M74/NGC 628 & 螺旋星系 & - & \(\ang{;10;30}\times\ang{;9;30}\) & \(1^\text{h }36^\text{m }42^\text{s}\) & \(+\ang{15;47;01}\) \\ %RGB/LRGBH\(\alpha\) & - \\ 
\end{longtable}
\end{center}


\subsection{摩羯座 Capricornus}
\begin{enumerate}
    \item 黃道帶最小的星座
    \item 摩羯一詞來源於摩伽羅，即印度神話中的海獸，而身體和尾部呈魚形，和希臘神話中Capricorn的形象吻合
\end{enumerate}
摩羯座中的深空天體：
\newline

\begin{center}
\begin{longtable}{c | c | c | c | c | c}
    % \hline
    編號 & 種類 & 別名 & 視尺寸 & RA & Dec \\
    \hline
    \endfirsthead

    % \hline
    % 編號 & 種類 & 別名 & 視尺寸 & 拍攝技巧 & 曝光時間 \\
    % \hline
    % \endhead

    % \hline
    % \multicolumn{6}{r}{接下頁} \\
    \endfoot

    % \hline
    \endlastfoot

    M30/NGC 7099 & 球狀星團 & - & \(\ang{;12;}\) & \(21^\text{h }40^\text{m }22^\text{s}\) & \(-\ang{23;10;45}\) \\ %RGB & - \\ 
\end{longtable}
\end{center}


\subsection{顯微鏡座 Microscopium}
\begin{enumerate}
    \item 最亮星：璃瑜增一（顯微鏡座\(\gamma\)）黃色巨星，視星等4.68
    \item 曾被認為是人馬座的後腿
\end{enumerate}
顯微鏡座中的深空天體：
\newline

\begin{center}
\begin{longtable}{c | c | c | c | c | c}
    % \hline
    編號 & 種類 & 別名 & 視尺寸 & RA & Dec \\
    \hline
    \endfirsthead

    % \hline
    % 編號 & 種類 & 別名 & 視尺寸 & 拍攝技巧 & 曝光時間 \\
    % \hline
    % \endhead

    % \hline
    % \multicolumn{6}{r}{接下頁} \\
    \endfoot

    % \hline
    \endlastfoot

    NGC 6925 & 漩渦星系 & - & \(\ang{;3;6}\times\ang{;1;6}\) & \(20^\text{h }34^\text{m }21^\text{s}\) & \(-\ang{31;58;50}\) \\ %RGB & - \\ 
\end{longtable}
\end{center}

\subsection{天鶴座 Grus}
\begin{enumerate}
    \item 最亮星：鶴一（天鶴座\(\alpha\)），視星等1.7等，是顆藍白色的恆星
    \item 最初被認為是鄰近星座南魚座的一部分
    \item 名稱來自拉丁語的鶴
\end{enumerate}
天鶴座中的深空天體：
\newline

\begin{center}
\begin{longtable}{c | c | c | c | c | c}
    % \hline
    編號 & 種類 & 別名 & 視尺寸 & RA & Dec \\
    \hline
    \endfirsthead

    % \hline
    % 編號 & 種類 & 別名 & 視尺寸 & 拍攝技巧 & 曝光時間 \\
    % \hline
    % \endhead

    % \hline
    % \multicolumn{6}{r}{接下頁} \\
    \endfoot

    % \hline
    \endlastfoot

    \makecell{NGC 7552\\NGC 7582\\NGC 7590\\NGC 7599} & 星系 & 天鶴座四重奏 & \(\ang{1;10;}\times\ang{;30;}\) & \(23^\text{h }18^\text{m }22^\text{s}\) & \(-\ang{42;22;28}\) \\ %RGB & - \\ 
\end{longtable}
\end{center}

\subsection{杜鵑座 Tucana}
\begin{enumerate}
    \item 最亮星：鳥喙一，視星等2.86，代表大嘴鳥頭部
    \item 杜鵑座\(\beta\)星：4 等及 5 等兩顆恆星組成的雙星，主星再區分成 4.4 等和 4.5 等兩顆恆星
    \item 杜鵑座\(\kappa\)星：雙星，5 等及 8 等
    \item 拉丁語原名指南美洲大嘴鳥
    \item 在天上另外兩隻鳥（天鶴座及鳳凰座）N的南邊
\end{enumerate}
杜鵑座中的深空天體：
\newline

\begin{center}
\begin{longtable}{c | c | c | c | c | c}
    % \hline
    編號 & 種類 & 別名 & 視尺寸 & RA & Dec \\
    \hline
    \endfirsthead

    % \hline
    % 編號 & 種類 & 別名 & 視尺寸 & 拍攝技巧 & 曝光時間 \\
    % \hline
    % \endhead

    % \hline
    % \multicolumn{6}{r}{接下頁} \\
    \endfoot

    % \hline
    \endlastfoot

    C106/NGC 104 & 球狀星團 & 杜鵑座47（47 Tuc） & \(\ang{;43;48}\) & \(0^\text{h }24^\text{m }05^\text{s}\) & \(-\ang{72;04;53}\) \\ %RGB & - \\ 
    NGC 362 & 球狀星團 & - & \(\ang{;12;54}\) & \(1^\text{h }03^\text{m }14^\text{s}\) & \(-\ang{70;50;53}\) \\ %RGB & - \\
    SMC & 矮星系 & 小麥哲倫星系 & \(\ang{5;50}\times\ang{3;5}\) & \(0^\text{h }52^\text{m }37^\text{s}\) & \(-\ang{72;48;21}\) \\ %RGB/RGBHO & > 4 hr \\

\end{longtable}
\end{center}

